% Ambiente de quadro, para informação textual com bordas fechadas (NBR 14724,
% que remete às normas de apresentação tabular do IBGE). Tabela fica para
% dado numérico. Definição baseada no exemplo do abnTeX2, sobre a classe memoir.
\newcommand{\quadroname}{Quadro}
\newcommand{\listofquadrosname}{Lista de quadros}
\newfloat[chapter]{quadro}{loq}{\quadroname}
\newlistof{listofquadros}{loq}{\listofquadrosname}
\newlistentry{quadro}{loq}{0}
\counterwithout{quadro}{chapter}
\renewcommand{\cftquadroname}{\quadroname\space}
\renewcommand*{\cftquadroaftersnum}{\hfill\textendash\hfill}

% Quadros longos (xltabular/longtable) usam o contador e a lista de quadros
\usepackage{ltcaption}
% Linha parcial que não abre ponto de quebra (\cline abre) entre as duas linhas de um requisito
\usepackage{hhline}

% Quadros longos que passam de uma página, como na seção 8.3 das Normas de
% apresentação tabular do IBGE (1993): o título e o cabeçalho se repetem, e
% cada página indica (continua), (continuação) ou (conclusão). O cabeçalho do
% longtable é uma caixa montada uma vez, então a indicação é escrita na saída
% de cada página, na posição reservada no cabeçalho, com as páginas de início
% e de fim gravadas no .aux. Precisa de duas compilações.
%   \quadrolongo{nome}            registra o quadro (antes dele)
%   \quadromarcaposicao{nome-1}   lugar da indicação no primeiro cabeçalho
%   \quadromarcaposicao{nome-c}   lugar da indicação no cabeçalho repetido
%   \quadromarcapagina{nome-ini}  página de início (primeiro cabeçalho)
%   \quadromarcapagina{nome-fim}  página de fim (último rodapé)
\makeatletter
\newcommand*{\quadro@pag}[2]{\expandafter\gdef\csname quadro@pag@#1\endcsname{#2}}
\newcommand*{\quadro@pos}[3]{\expandafter\gdef\csname quadro@pos@#1\endcsname{{#2}{#3}}}
\newcommand*{\quadromarcapagina}[1]{\write\@auxout{\string\quadro@pag{#1}{\the\c@page}}}
\newcommand*{\quadromarcaposicao}[1]{%
	\pdfsavepos\write\@auxout{\string\quadro@pos{#1}{\the\pdflastxpos}{\the\pdflastypos}}}
\def\quadro@longos{}
\newcommand*{\quadrolongo}[1]{\g@addto@macro\quadro@longos{\quadro@indicacao{#1}}}
\newcommand*{\quadro@escreve}[2]{%
	\@ifundefined{quadro@pos@#1}{}{%
		\expandafter\expandafter\expandafter\quadro@escreve@\csname quadro@pos@#1\endcsname{#2}}}
\newcommand*{\quadro@escreve@}[3]{%
	\put(0,0){\rlap{\kern#1sp\raise\dimexpr#2sp-\paperheight\relax\hbox{\llap{\normalfont\footnotesize #3}}}}}
\newcommand*{\quadro@indicacao}[1]{%
	\@ifundefined{quadro@pag@#1-ini}{}{\@ifundefined{quadro@pag@#1-fim}{}{%
		\edef\quadro@ini{\csname quadro@pag@#1-ini\endcsname}%
		\edef\quadro@fim{\csname quadro@pag@#1-fim\endcsname}%
		\ifnum\quadro@ini<\quadro@fim\relax
			\ifnum\c@page=\quadro@ini\relax \quadro@escreve{#1-1}{(continua)}\fi
			\ifnum\c@page>\quadro@ini\relax
				\ifnum\c@page<\quadro@fim\relax \quadro@escreve{#1-c}{(continuação)}\fi
				\ifnum\c@page=\quadro@fim\relax \quadro@escreve{#1-c}{(conclusão)}\fi
			\fi
		\fi}}}
\AddToHook{shipout/foreground}{\quadro@longos}
\makeatother
